\documentclass[12pt,a4paper]{article}

\usepackage[spanish,es-noshorthands]{babel}
\usepackage[margin=2.2cm]{geometry}
\usepackage{graphicx}
\usepackage{tikz}
\usepackage{xcolor}
\usepackage{enumitem}
\usepackage{hyperref}
\usepackage{fancyhdr}

\definecolor{rojo}{HTML}{C74221}
\definecolor{verde}{HTML}{1B7A3D}
\definecolor{gris}{HTML}{F2F2F2}
\definecolor{atlasCoverBg}{HTML}{20242E}
\definecolor{atlasMuted}{HTML}{9AA3B2}
\definecolor{atlasMarino}{HTML}{010423}
\definecolor{atlasDorado}{HTML}{F7C65F}

\pagestyle{fancy}
\fancyhf{}
\fancyhead[L]{\small Manual POS --- Textiles y Confecciones Atlas}
\fancyfoot[C]{\small P\'agina \thepage}
\setlength{\headheight}{14pt}

% Foto con recuadros numerados en coordenadas relativas (0..1).
% Uso:
% \foto{imagenes/02.png}{%
%   \recuadro{0.42}{0.10}{0.60}{0.25}{1}%
% }
\newcommand{\recuadro}[5]{%
  \draw[white,line width=3.5pt] (#1,#2) rectangle (#3,#4);%
  \draw[rojo,ultra thick] (#1,#2) rectangle (#3,#4);%
  \node[fill=rojo,text=white,font=\bfseries\small,inner sep=3pt,rounded corners=2pt] at (#3+0.030,#4-0.005) {#5};%
}
\newcommand{\foto}[2]{%
  \begin{center}
  \begin{tikzpicture}
    \node[anchor=south west,inner sep=0] (img) at (0,0)
      {\includegraphics[width=0.96\linewidth]{#1}};
    \begin{scope}[x={(img.south east)},y={(img.north west)}]
      #2
    \end{scope}
  \end{tikzpicture}
  \end{center}
}

\newcommand{\paso}[1]{\par\medskip\noindent\textbf{Paso #1.} }
\newcommand{\punto}[1]{\par\medskip\noindent\textbf{Punto #1.} }
\newcommand{\aviso}[1]{%
  \par\smallskip\noindent
  \fcolorbox{rojo}{gris}{\parbox{\dimexpr\linewidth-2\fboxsep-2\fboxrule}{#1}}
  \par\smallskip
}
\newcommand{\consejo}[1]{%
  \par\smallskip\noindent
  \fcolorbox{verde}{gris}{\parbox{\dimexpr\linewidth-2\fboxsep-2\fboxrule}{#1}}
  \par\smallskip
}

\begin{document}

% ---------------- PORTADA ----------------
\begin{titlepage}
\thispagestyle{empty}

% Banda superior azul marino
\noindent\colorbox{atlasMarino}{\parbox{\dimexpr\linewidth-2\fboxsep}{%
  \color{white}\small\bfseries
  MANUAL DE USUARIO%
  \hfill TEXTILES Y CONFECCIONES ATLAS E.I.R.L.%
}}
\vspace{1.4cm}

\begin{center}
  \includegraphics[width=6.5cm]{imagenes/logo-atlas.png}
\end{center}
\vspace{0.7cm}

\begin{center}
  {\color{atlasMarino}\large\bfseries\textsc{Manual de usuario}}\\[4pt]
  {\color{rojo}\rule{9cm}{2.2pt}}\\[2pt]
  {\color{atlasDorado}\rule{9cm}{1pt}}\\[12pt]
  {\color{atlasMarino}\LARGE\bfseries Punto de Venta (POS)}\\[6pt]
  {\color{atlasMarino}\Large\bfseries Venta con boleta o factura electr\'onica}\\[10pt]
  {\large\color{black!70} Gu\'ia operativa ilustrada para el personal de tienda}\\[4pt]
  {\normalsize\color{black!70} Sistema Odoo \textperiodcentered\ Tienda de ropa}
\end{center}
\vspace{0.8cm}

% Ficha del documento
\noindent\fcolorbox{atlasMarino}{gris}{\parbox{\dimexpr\linewidth-2\fboxrule-2\fboxsep}{%
  \vspace{2pt}
  \begin{tabular}{@{\hspace{6pt}}l@{\hspace{12pt}}l@{}}
    \color{atlasMarino}\textbf{Empresa:} &
      Textiles y Confecciones Atlas E.I.R.L. \textperiodcentered\ RUC 20558074427 \\[4pt]
    \color{atlasMarino}\textbf{Alcance:} &
      Venta mostrador con boleta o factura electr\'onica (SUNAT) \\[4pt]
    \color{atlasMarino}\textbf{Dirigido a:} &
      Vendedoras y vendedores de tienda \\[4pt]
    \color{atlasMarino}\textbf{Versi\'on:} &
      1.0 \textperiodcentered\ Octubre 2026 \\
  \end{tabular}%
  \vspace{2pt}
}}
\vfill

% Banda inferior azul marino
\noindent\colorbox{atlasMarino}{\parbox{\dimexpr\linewidth-2\fboxsep}{%
  \centering\color{white}\small
  Documento de uso interno \textperiodcentered\ Prohibida su reproducci\'on sin autorizaci\'on%
}}
\end{titlepage}

\tableofcontents
\newpage

% ---------------- CAP 1 ----------------
\clearpage
\section{Antes de empezar: lo m\'inimo que debe saber}

  \punto{1} Encienda la computadora, abra el navegador de internet (Chrome o Edge)
  y entre a la p\'agina de Odoo de la empresa con su usuario y contrase\~na.
  Si no tiene usuario, p\'idaselo a su encargado. No comparta su contrase\~na.

  \punto{2} Este manual usa siempre las mismas palabras:
  \begin{itemize}
    \item \textbf{Hacer clic} = presionar una vez el bot\'on izquierdo del mouse
    sobre lo que se indica.
    \item \textbf{Pantalla} = lo que usted ve en el monitor.
    \item \textbf{Bot\'on} = un recuadro con texto que se puede presionar.
  \end{itemize}

  \punto{3} \textbf{Idea m\'as importante del manual:}
  en la pantalla de cobro hay un bot\'on que dice \textbf{Recibo/Factura}.
  \begin{itemize}
    \item Si ese bot\'on tiene el cuadrito \textbf{marcado con check},
    la venta \textbf{S\'I se env\'ia a la SUNAT} (boleta o factura electr\'onica).
    \item Si el cuadrito est\'a \textbf{vac\'io}, la venta \textbf{NO se env\'ia a la SUNAT}.
  \end{itemize}
  La mayor\'ia de las ventas de la tienda \textbf{s\'i} deben enviarse a la SUNAT.
  No quite el check salvo que su encargado se lo pida.

\aviso{\textbf{Boleta o factura, \'cu\'al sale?} Usted no elige a mano.
El sistema decide solo seg\'un el cliente: si el cliente tiene \textbf{DNI},
sale \textbf{boleta}; si tiene \textbf{RUC}, sale \textbf{factura}.
Por eso es muy importante elegir bien al cliente (cap\'itulo 5).}

% ---------------- CAP 2 ----------------
\clearpage
\section{Entrar al Punto de venta}

\paso{1} Despu\'es de ingresar a Odoo ver\'a el men\'u de aplicaciones
(cuadritos con dibujos). Haga clic en el cuadrito
\textbf{Punto de venta} (dibujo de una tiendita).

\foto{imagenes/00-home-odoo.png}{%
  \recuadro{0.2049}{0.7944}{0.3032}{0.9222}{1}%
}

\paso{2} Ver\'a la tarjeta de la tienda (\textbf{Tienda de ropa}).
Haga clic en el bot\'on morado \textbf{Seguir vendiendo}.

\foto{imagenes/01-pos-dashboard.png}{%
  \recuadro{0.0229}{0.7642}{0.1196}{0.8009}{1}%
}

\consejo{Si en lugar de \textbf{Seguir vendiendo} ve un bot\'on que dice
\textbf{Abrir caja} o \textbf{Nueva sesi\'on}, avise a su encargado:
la caja del d\'ia todav\'ia no est\'a abierta y solo \'el debe abrirla.}

% ---------------- CAP 3 ----------------
\clearpage
\section{La pantalla de venta: qu\'e es cada cosa}

Esta es la pantalla donde usted trabajar\'a todo el d\'ia.
M\'irela con calma:

\foto{imagenes/02-pos-pantalla-principal.png}{%
  \recuadro{0.0056}{0.9361}{0.0730}{0.9906}{1}%
  \recuadro{0.7174}{0.9476}{0.8300}{0.9790}{2}%
  \recuadro{0.5908}{0.6080}{0.7217}{0.8189}{3}%
  \recuadro{0.0056}{0.0178}{0.1192}{0.0722}{4}%
}

  Los n\'umeros rojos de la foto significan:

  \textbf{1.} Pesta\~nas \textbf{Registrar} y \textbf{\'Ordenes}:
  \textbf{Registrar} es la venta actual (la que est\'a atendiendo).
  \textbf{\'Ordenes} es la lista de ventas ya terminadas. Trabaje siempre en
  \textbf{Registrar}.

  \textbf{2.} \textbf{Buscar productos:} aqu\'i escribe el nombre del producto
  con el teclado (por ejemplo: \texttt{polo}) y el sistema lo filtra solo.

  \textbf{3.} \textbf{Productos:} cada tarjeta con foto es un producto.
  Para agregarlo a la venta, haga clic sobre su tarjeta (cap\'itulo 4).

  \textbf{4.} \textbf{Consumidor Final:} es el cliente de la venta actual.
  Al inicio siempre dice \textit{Consumidor Final}. Debe cambiarlo por
  el cliente real (cap\'itulo 5). Abajo est\'a el bot\'on morado de
  \textbf{Pago} para cobrar (cap\'itulo 7).

% ---------------- CAP 4 ----------------
\clearpage
\section{Agregar productos a la venta}

\paso{1} Busque el producto (recuadro de b\'usqueda) o ub\'iquelo con la vista
en sus tarjetas.

\paso{2} Haga \textbf{un solo clic} sobre la tarjeta del producto.
Cada clic agrega \textbf{una unidad}. Si el cliente lleva 3 polos iguales,
haga clic 3 veces en el mismo polo.

\paso{3} Revise el lado izquierdo de la pantalla: aparece la l\'inea del producto,
la cantidad, los impuestos y el \textbf{Total} en soles. Ejemplo de la foto:
\textit{1 Producto Prueba 01 --- S/ 1.00}.

\foto{imagenes/03-producto-agregado.png}{%
  \recuadro{0.0000}{0.8861}{0.3118}{0.9256}{1}%
  \recuadro{0.5908}{0.6080}{0.7217}{0.8189}{2}%
  \recuadro{0.0056}{0.4107}{0.3062}{0.4360}{3}%
  \recuadro{0.0056}{0.0089}{0.3062}{0.0867}{4}%
}

  Los n\'umeros rojos de la foto significan:

  \textbf{1.} La l\'inea de la venta: cantidad, nombre y precio.
  Si se equivoc\'o de producto, haga clic sobre esa l\'inea y luego en el
  bot\'on rojo con la X para borrarla.

  \textbf{2.} El numerito negro sobre la tarjeta confirma cu\'antas unidades
  lleva de ese producto.

  \textbf{3.} El \textbf{Total}. Debe coincidir con lo que el cliente
  va a pagar.

  \textbf{4.} El bot\'on \textbf{Pago}. Cuando la lista de productos est\'e
  completa, haga clic ah\'i para cobrar. Todav\'ia NO ha cobrado nada:
  todo se puede corregir.

\aviso{\textbf{No haga clic dos veces seguidas sin querer.}
Si el mouse hace doble clic, se agregan 2 unidades. Mire siempre
el numerito de la tarjeta y el Total antes de seguir.}

% ---------------- CAP 5 ----------------
\clearpage
\section{Elegir al cliente: el paso que decide boleta o factura}

\paso{1} Debajo del Total, haga clic en el bot\'on que dice
\textbf{Consumidor Final}.

\paso{2} Se abre la ventana \textbf{Elige un cliente} con la lista de clientes.
Tiene tres caminos:
\begin{itemize}
  \item \textbf{El cliente ya est\'a en la lista:} haga clic sobre su nombre.
  \item \textbf{Es un cliente nuevo:} haga clic en el bot\'on morado
  \textbf{Crear} y llene sus datos (ver recuadro de abajo).
  \item \textbf{Es una venta r\'apida sin datos (monto peque\~no):}
  puede dejar \textbf{Consumidor Final} y seguir. Saldr\'a una boleta simple.
\end{itemize}

\foto{imagenes/04-lista-clientes.png}{%
  \recuadro{0.1715}{0.8956}{0.2191}{0.9500}{1}%
  \recuadro{0.6384}{0.9071}{0.7932}{0.9385}{2}%
  \recuadro{0.1608}{0.5883}{0.8392}{0.6950}{3}%
}

  Los n\'umeros rojos de la foto significan:

  \textbf{1.} \textbf{Crear:} solo para clientes nuevos. No cree dos veces
  al mismo cliente: primero b\'usquelo arriba con su nombre o documento.

  \textbf{2.} \textbf{Buscar clientes:} escriba el nombre, DNI o RUC y presione
  Enter.

  \textbf{3.} \textbf{La fila del cliente:} haga clic sobre el nombre para elegirlo.
  F\'ijese en el documento: los clientes con \textbf{RUC} (empresas) salen con
  su direcci\'on completa; los de \textbf{DNI} (personas) salen m\'as simples.

\subsection{Si el cliente es nuevo: c\'omo llenar la ficha}

Al hacer clic en \textbf{Crear} se abre la ficha \textbf{Crear contacto}.
Ll\'enela as\'i, de arriba hacia abajo:

\foto{imagenes/08-crear-cliente.png}{%
  \recuadro{0.2840}{0.7533}{0.8174}{0.8011}{1}%
  \recuadro{0.1771}{0.4378}{0.4549}{0.4911}{2}%
  \recuadro{0.1715}{0.1327}{0.2216}{0.1694}{3}%
}

  \paso{1} \textbf{Nombre:} escriba el nombre completo de la persona o empresa,
  tal como figura en su DNI o RUC. Sin apodos.
  \paso{2} \textbf{N\'umero de Identificaci\'on:} aqu\'i est\'a el secreto de
  la boleta y la factura:
  \begin{itemize}
    \item Abra la lista y elija \textbf{DNI} si es una persona
    (8 n\'umeros). La venta saldr\'a como \textbf{BOLETA}.
    \item Elija \textbf{RUC} si es una empresa (11 n\'umeros).
    La venta saldr\'a como \textbf{FACTURA}.
  \end{itemize}
  Escriba el n\'umero con cuidado, d\'igito por d\'igito. Un n\'umero mal
  escrito hace que la SUNAT rechace el comprobante.
  \paso{3} Revise todo y haga clic en \textbf{Guardar}.
  Si se arrepiente, haga clic en \textbf{Descartar}: no se guarda nada.

\aviso{\textbf{Regla de oro:} DNI = boleta, RUC = factura.
Si el cliente pide factura pero solo le da su DNI, expl\'iquele con amabilidad
que sin RUC el sistema solo puede emitir boleta.}

\subsection{Cliente ya elegido}

Al elegir al cliente, el bot\'on de abajo cambia y muestra su nombre.
En la foto se eligi\'o a \textbf{COLEGIO M\'EDICO DEL PER\'U} (tiene RUC,
entonces esta venta saldr\'a como \textbf{factura}).

\foto{imagenes/05-cliente-seleccionado.png}{%
  \recuadro{0.0056}{0.3374}{0.1793}{0.3918}{1}%
  \recuadro{0.0056}{0.0089}{0.3062}{0.0867}{2}%
}

  \paso{1} Compruebe que el nombre del cliente es el correcto.
  Si se equivoc\'o, haga clic otra vez sobre el nombre y elija otro.
  \paso{2} Todo listo para cobrar: haga clic en \textbf{Pago}.

% ---------------- CAP 6 ----------------
\clearpage
\section{Cobrar: la pantalla de pago y el bot\'on Recibo/Factura}

Al hacer clic en \textbf{Pago} aparece la pantalla de cobro.
Tiene tres zonas:

\foto{imagenes/06-pantalla-pago.png}{%
  \recuadro{0.0056}{0.8456}{0.3069}{0.9167}{1}%
  \recuadro{0.5549}{0.3578}{0.7576}{0.5667}{2}%
  \recuadro{0.1590}{0.3363}{0.3069}{0.3907}{3}%
}

  \paso{1} \textbf{Medio de pago:} c\'omo paga el cliente
  (Efectivo, Tarjeta, YAPE, PLIN o Cuenta de cliente).
  Haga clic en el que corresponda. Este manual explica el cobro en
  \textbf{Efectivo}, que es el m\'as com\'un.
  \paso{2} \textbf{El monto grande:} es el total a cobrar (ejemplo: S/ 1.00).
  Debe coincidir con lo que usted le dice al cliente.
  \paso{3} \textbf{Recibo/Factura con check:} significa que al cobrar,
  el sistema \textbf{crear\'a el comprobante y lo enviar\'a a la SUNAT}.
  D\'ejelo marcado. Si est\'a vac\'io (foto siguiente), la venta
  \textbf{NO} genera comprobante para la SUNAT.

\subsection{Con check y sin check: la diferencia}

Compare las dos fotos. En la primera el cuadrito tiene check:
\textbf{esta venta S\'I va a la SUNAT}. En la segunda el cuadrito est\'a vac\'io:
\textbf{esta venta NO va a la SUNAT} (solo queda registrada en la tienda).

\foto{imagenes/07-tipo-comprobante.png}{%
  \recuadro{0.1590}{0.3363}{0.3069}{0.3907}{1}%
}

\aviso{\textbf{Si toca el bot\'on Recibo/Factura, cambia de un estado al otro.}
No lo toque por curiosidad en una venta real: si lo deja sin check,
el cliente se ir\'a sin su boleta o factura electr\'onica y habr\'a que
arreglarlo despu\'es con su encargado y el contador.}

\consejo{Antes de seguir: mire las tres cosas en orden:
1) medio de pago correcto, 2) monto correcto, 3) Recibo/Factura con check.
Solo entonces contin\'ue.}

% ---------------- CAP 7 ----------------
\clearpage
\section{Cobrar en efectivo paso a paso}

\paso{1} En la pantalla de pago, haga clic en \textbf{Efectivo}.

\paso{2} El sistema muestra \textbf{Procesando\dots} con el monto grande.
Espere sin tocar nada: est\'a guardando la venta y preparando el comprobante.

\foto{imagenes/09-pago-efectivo.png}{%
  \recuadro{0.2778}{0.6580}{0.7222}{0.7780}{1}%
}

\paso{3} Cuando termine (desaparece el recuadro 1 y aparecen los botones \textbf{Imprimir},
\textbf{Enviar recibo} y \textbf{Continuar}:
\begin{itemize}
  \item \textbf{Imprimir:} imprime el ticket para entregar al cliente.
  \item \textbf{Enviar recibo:} se lo env\'ia al correo del cliente (si tiene).
  \item \textbf{Continuar:} vuelve a la pantalla de venta para atender
  al siguiente cliente. La venta anterior ya qued\'o guardada.
\end{itemize}

\paso{4} Haga clic en \textbf{Continuar} y verifique que la pantalla qued\'o
limpia (sin productos) antes de atender al siguiente cliente.

\aviso{\textbf{Un clic, una sola vez.} Cuando est\'a en \textbf{Procesando\dots},
NO haga clic de nuevo ni cierre la ventana: puede duplicar el cobro.
Espere a que aparezcan los botones Imprimir y Continuar.}

% ---------------- CAP 8 ----------------
\clearpage
\section{Errores frecuentes y qu\'e hacer}

\begin{enumerate}[leftmargin=*]
  \item \textbf{Agregu\'e un producto de m\'as (o el equivocado).} \\
  Antes de cobrar: haga clic en la l\'inea del producto (lado izquierdo)
  y b\'orrela con el bot\'on rojo X. Despu\'es de cobrar ya NO se puede
  borrar: avise a su encargado (se corrige con nota de cr\'edito).

  \item \textbf{Eleg\'i al cliente equivocado.} \\
  Antes de cobrar: haga clic otra vez sobre el nombre del cliente y elija
  el correcto. Despu\'es de cobrar ya NO se puede cambiar.

  \item \textbf{Olvid\'e marcar Recibo/Factura y cobr\'e sin comprobante.} \\
  Avise de inmediato a su encargado. Esa venta qued\'o sin enviar a la SUNAT
  y debe regularizarse el mismo d\'ia.

  \item \textbf{La pantalla se qued\'o en Procesando\dots mucho tiempo.} \\
  Espere un minuto sin tocar nada. Si no avanza, avise a su encargado.
  No intente cobrar la misma venta en otra computadora: podr\'ia duplicarse.

  \item \textbf{El cliente dice que su boleta/factura ``no aparece''.} \\
  Los comprobantes tienen plazos de la SUNAT: la boleta individual hasta
  5 d\'ias y el resumen hasta 7 d\'ias; la factura hasta 7 d\'ias.
  Si el env\'io se demor\'o, el sistema muestra un error y el comprobante
  queda pendiente. Avise a su encargado para que lo regularice cuanto antes.
\end{enumerate}

\clearpage
\section*{Resumen de una venta normal (para pegar en caja)}

\begin{enumerate}
  \item Entrar a \textbf{Punto de venta} $\rightarrow$ \textbf{Seguir vendiendo}.
  \item Clic en cada producto que lleva el cliente. Revisar el Total.
  \item Clic en \textbf{Consumidor Final} $\rightarrow$ buscar o crear al cliente
  (DNI = boleta, RUC = factura).
  \item Clic en \textbf{Pago}. Revisar: medio de pago, monto y
  \textbf{Recibo/Factura con check}.
  \item Clic en el medio de pago (ej. \textbf{Efectivo}). Esperar sin tocar nada.
  \item \textbf{Imprimir} el ticket, entregar al cliente y hacer clic en
  \textbf{Continuar}.
\end{enumerate}

\end{document}
